%%%%%%%%%%%%%%%%%%%%%%%%%%%%%%%%%%%%%%%%%%%%%%%%%%%%%%%%%%%%%%%%%%%%%%%%%%%%%%
% Short Sectioned Assignment LaTeX Template Version 1.0 (5/5/12)
% This template has been downloaded from: http://www.LaTeXTemplates.com
% Original author:  Frits Wenneker (http://www.howtotex.com)
% License: CC BY-NC-SA 3.0 (http://creativecommons.org/licenses/by-nc-sa/3.0/)
%%%%%%%%%%%%%%%%%%%%%%%%%%%%%%%%%%%%%%%%%%%%%%%%%%%%%%%%%%%%%%%%%%%%%%%%%%%%%%


\documentclass[11pt, a4paper]{book}

% PAQUETES ------------------------------------------------------------------------------

\usepackage[T1]{fontenc}
\usepackage[utf8]{inputenc}
% Selecciona el español para palabras introducidas automáticamente, p.ej. "septiembre" en la fecha y especifica que se use la palabra Tabla en vez de Cuadro
\usepackage[spanish,es-tabla,es-nodecimaldot]{babel}
\usepackage{amsmath}
\usepackage{amsfonts}
\usepackage{amssymb}
% Manejo de símbolos de cita múltiples
\usepackage{csquotes}
% Figuras e imágenes
\usepackage{graphics, graphicx, float}
% Permite dar estilo a los subtítulos de las figuras (leyendas)
\usepackage{caption}
% Permite sub-leyendas en sub-figuras (figuras dentro de figuras)
\usepackage{subcaption}
% Permite múltiples filas en tablas
\usepackage{multirow}
% Elimina numeración de paginas en blanco
\usepackage{emptypage}
% Permite continuar con la numeración entre listas distintas
\usepackage{enumitem}
% Introduce la bibliografía en el índice
\usepackage[nottoc]{tocbibind}
% Permite continuar con la numeración entre listas distintas
\usepackage{fourier} % Use the Adobe Utopia font for the document
\usepackage{url} % ,href} %para incluir URLs e hipervínculos dentro del texto (aunque hay que instalar href)
% Para incluir el símbolo del euro
\usepackage[gen]{eurosym} 
\usepackage{cite} %para incluir citas del archivo <nombre>.bib
\usepackage{enumerate}
% Permite añadir comentarios multilinea que no se mostrarán en el PDF final
\usepackage{comment}
\usepackage{tabularx}
\usepackage{booktabs}
\usepackage[table,xcdraw]{xcolor}
% Marca en rojo los elementos pendientes (TO-DOs)
\newcommand{\todo}[1]{\textcolor{orange}{TODO: #1}}
% Permite que el contador de footnotes no se reinicie en cada capítulo
\usepackage{chngcntr}
\counterwithout{footnote}{chapter}

% ESTILOS DE PAGINA ------------------------------------------------------------------
\renewcommand{\familydefault}{\sfdefault}
% Modificar encabezamientos y pie de páginas
\usepackage{fancyhdr}
\pagestyle{fancyplain} % Makes all pages in the document conform to the custom headers and footers
\fancyhead[L]{} % Empty left header
\fancyhead[C]{} % Empty center header
% \fancyhead[R]{Francisco David Castejón Soto} % Nombre del autor en el encabezado derecho
\fancyfoot[L]{} % Empty left footer
\fancyfoot[C]{} % Empty center footer
\fancyfoot[R]{\thepage} % Page numbering for right footer
\renewcommand{\headrulewidth}{0pt} % Remove header underlines
\renewcommand{\footrulewidth}{0pt} % Remove footer underlines
\setlength{\headheight}{13.6pt} % Customize the height of the header

\usepackage{titlesec, blindtext, color}
\definecolor{gray75}{gray}{0.75}
\newcommand{\hsp}{\hspace{20pt}}
\titleformat{\chapter}[hang]{\Huge\bfseries}{\thechapter\hsp\textcolor{gray75}{|}\hsp}{0pt}{\Huge\bfseries}
\setcounter{secnumdepth}{4}
\usepackage[Lenny]{fncychap}

% For code snippets
\usepackage{minted}
\usemintedstyle{colorful}
% Control de enlaces hipertexto en el documento. Este paquete debe cargarse el último.
\usepackage{hyperref}
\hypersetup{
	colorlinks=true,	% false: boxed links; true: colored links
	linkcolor=black,	% color of internal links
	urlcolor=cyan		% color of external links
}

% Información del documento ---------------------------------------------------------
\author{Francisco David Castejón Soto}
\title{Entrenamiento de una IA mediante aprendizaje por refuerzo para un juego hecho en Unreal Engine}
\date{\today}

% DOCUMENTO ----------------------------------------------------------------------------
\begin{document}

% Páginas preliminares sin numeración
\frontmatter

\begin{titlepage}
	\thispagestyle{empty}

	\centering
	\includegraphics[width=0.9\textwidth]{logos/logo_ugr.jpg}
	\vspace{1.0cm}

	{\rmfamily\textsc{\Large MÁSTER UNIVERSITARIO}}\\[0.2cm]
	{\rmfamily\textsc{\large en Ciencia de Datos e Ingeniería de Computadores}}
	\vspace{1.5cm}

	\vfill

	{\Huge\bfseries Trabajo de Fin de Máster}\\[0.5cm]
	\rule{\textwidth}{3pt}\\[0.5cm]
	{\Large\bfseries Desarrollo y Optimización de un Agente Inteligente \\
	para Scripts of Tribute mediante Algoritmos Evolutivos}
	\vspace{1.5cm}

	\vfill

	\begin{tabular}{@{}c@{}}
		\textbf{\large Autor}         \\[0.3cm]
		Francisco David Castejón Soto \\[1cm]
		\textbf{\large Director}      \\[0.3cm]
		Dr. Pablo García Sánchez
	\end{tabular}

	\vfill

	\includegraphics[width=0.4\textwidth]{logos/etsiit_logo.png}\\[0.3cm]
	{\rmfamily\textsc{\footnotesize Escuela Técnica Superior de Ingenierías Informática y de Telecomunicación}}\\
	\rule{0.1\textwidth}{0.5pt}\\[0.3cm]
	{\large Granada, 1 Julio de 2025}

\end{titlepage}
\chapter*{Prefacio}

Para la realización de este documento, se ha utilizado la plantilla \LaTeX \cite{guervos_jjplantilla-tfg-etsiit_2024} específica para la ETSIIT.
\chapter*{Agradecimientos} \label{ch:agradecimientos}

A mi familia por el apoyo moral y económico que me han brindado de forma incondicional durante estos años fuera de casa. Y a mi pareja, por mostrarme más caminos de los que nunca hubiera pensado que existieran para mí.

% Índices
\cleardoublepage
\tableofcontents

\cleardoublepage
\listoffigures

\cleardoublepage
\listoftables

% Contenido principal
\mainmatter

\part{Marco teórico: Cartas y Bots} \label{part:idea}
\chapter{Introducción} \label{chap:introduccion}

La inteligencia artificial (IA) en videojuegos existe en una intersección entre la ciencia computacional y el arte del entretenimiento, dando vida a mundos virtuales y creando oponentes dignos de nuestras mejores estrategias. En este contexto, la IA se refiere únicamente a un conjunto de algoritmos y técnicas diseñadas para cumplir una función muy específica dentro del videojuego. Ejemplos de esto son la IA que ``conduce'' los coches en un videojuego de carreras, los Pokémon contra los que el jugador se enfrenta en una batalla o el simple movimiento de una línea de píxeles (a modo de pala contrincante) en el videojuego Pong. Esta acepción específica contrasta con la idea más extendida de IA, que posee un conjunto de connotaciones más amplias y generalistas como ``la máquina que aprende'' o ``el estudio y construcción de sistemas que hacen <<lo correcto>>, en función de su objetivo'' \cite{russell_artificial_2020}. Independientemente de la definición que se utilice, un factor clave en el desarrollo de un videojuego es la creación de una IA que resulte entretenida para el jugador. De la misma forma que en un libro o una película es necesaria una trama que plantee un desafío, como la lucha contra el Imperio Galáctico en Star Wars; o un compañero de aventuras que ayude al protagonista, como Sam en El Señor de los Anillos. En un videojuego es necesaria una IA que controle esos elementos que interactúan con el jugador, ya sean enemigos, aliados o incluso el propio entorno. Desde los propios inicios de la industria del videojuego, donde se construían máquinas específicas para ejecutar un juego concreto, como es el caso de Nim en 1948 \cite{redheffer_machine_1948}, hasta los videojuegos modernos, donde se utilizan técnicas de aprendizaje automático complejas, como el aprendizaje por refuerzo \cite{gaudreau_game_2025}, la IA ha sido un componente esencial para crear experiencias de juego divertidas y desafiantes.

La intención de este trabajo es la de crear una IA contra la que jugar en un videojuego de estrategia. En concreto, la que controla a un agente autónomo (o simplemente ``bot'') en un simulador del videojuego de cartas \textit{Tales of Tribute}\footnote{Es importante la distinción entre ``juego'' y ``videojuego'' en este caso, ya que las cartas no son físicas y tangibles, sino que únicamente existen en el universo digital del videojuego.}. Para guiar las decisiones del bot, un algoritmo evolutivo se encarga de ajustar los pesos que se utilizan para calcular la puntuación de cada jugada posible, eligiendo vorazmente la jugada que maximiza dicha puntuación. De esta forma se ha conseguido crear un contrincante que juega de forma competitiva tanto contra el jugador humano como contra otros bots manejados por algoritmos de inteligencia artificial totalmente diferentes.
\chapter{Objetivos} \label{chap:objetivos}
% Definir lo que se busca conseguir con el TFM


\section{Objetivos principales} \label{sec:objetivo_principal}

Los objetivos principales de este Trabajo de Fin de Máster representan dos metas entrelazadas. Por un lado, ahondar en la investigación y entendimiento de un área tan extensa como las metaheurísticas, y en concreto los algoritmos evolutivos. La optimización de soluciones a problemas complejos mediante técnicas que tratan de incluir conocimiento específico, unido a la generación iterativa de otras soluciones parciales, es un campo de estudio con un gran número de aplicaciones prácticas. Aunque el teorema de ``No Free Lunch'' nos advierte de que no existe una única técnica que sea la mejor para todos los problemas y desde todas las perspectivas \cite{wolpert_no_1997}, la experiencia empírica ha demostrado que los algoritmos evolutivos son una herramienta que se puede aplicar a la gran mayoría de problemas de optimización \cite{torres-jimenez_applications_2014}.

Es precisamente esa versatilidad la que ha propiciado su uso durante el desarrollo del segundo objetivo principal de este proyecto: la creación de un bot para un videojuego de cartas. Este tipo de videojuegos cuentan con un inmenso número de variables, pues no solo se deben de tener en cuenta las cartas existentes en las manos de cada jugador y en la pila de cartas, sino también las mecánicas\footnote{Gran parte de lo que define un juego es sus mecánicas, sus reglas, lo que los jugadores deben llevar a cabo para ganar. La labor del creador de videojuegos es conseguir un conjunto de reglas que estén equilibradas y permitan disfrutar a los jugadores \cite{wikipedia_diseno_2025}.} intrínsecas del juego, como la vida, los recursos de compra, la posibilidad de recuperar cartas usadas, o el uso de otro tipo de habilidades especiales. Todo esto hace que la creación de un bot que juegue a alto nivel a un videojuego de este estilo sea un reto interesante, y que el uso de técnicas de optimización evolutiva sea una herramienta adecuada para la optimización de su rendimiento.

\section{Objetivos específicos} \label{sec:objetivos_especificos}

% Objetivos del formulario de petición de tema:
\begin{itemize}
	\item OG1: Analizar en profundidad las mecánicas del juego ``Tales of Tribute'', el entorno de desarrollo ``Scripts of Tribute'' y adquirir las competencias necesarias en el lenguaje de programación C\#.
	\item OG2: Diseñar e implementar un agente inteligente en C\# para ``Scripts of Tribute'', cuya toma de decisiones se base en una evaluación heurística del estado del juego mediante una función de fitness ponderada, incorporando conocimiento específico del dominio.
	\item OG3: Desarrollar un programa de optimización en Python para ajustar los pesos de la función de fitness del agente, implementando y comparando dos estrategias principales: algoritmos coevolutivos y entrenamiento supervisado contra agentes de referencia de ``Scripts of Tribute''.
	\item OG4: Desarrollar un conjunto de herramientas para la visualización y análisis de los datos generados durante el entrenamiento.
	\item OG5: Evaluar cuantitativamente el rendimiento del agente entrenado mediante las diferentes estrategias, utilizando métricas relevantes.
	\item OG6: Analizar comparativamente la efectividad y eficiencia de las estrategias de optimización implementadas, discutiendo sus ventajas y desventajas en el contexto específico de ``Scripts of Tribute''.
	\item OG7: Investigar y contextualizar el enfoque desarrollado frente a otras técnicas predominantes en competiciones similares de IA en juegos, analizando las razones de su éxito.
\end{itemize}
\chapter{Antecedentes y estado del arte} \label{chap:antecedentes}

\section{Inteligencia artificial en videojuegos comerciales} \label{sec:ia_videojuegos}

En esta sección se realiza un breve repaso a las técnicas de inteligencia artificial más utilizadas a lo largo de la historia de los videojuegos, así como algunos ejemplos del estado del arte en el ámbito de la industria comercial del videojuego. Los usos de estos algoritmos no solo abarcan su función más habitual de enemigo o aliado, sino que también abarcan otros aspectos como la búsqueda de caminos o la deformación realista de las mallas \footnote{Una ``mesh'', o malla en español, es una estructura tridimensional formada por una colección de vértices, aristas y caras que definen la forma de un objeto 3D dentro de un videojuego \cite{universidad_europea_que_2025}.} de los personajes animados.

\subsection{Funciones hash}

Una de las técnicas más sencillas y antiguas para controlar partes específicas de un videojuego son las funciones hash. Este tipo de funciones simplemente reciben un conjunto de parámetros como entrada y devuelven un valor o acción a realizar. Dada su simplicidad, generan comportamientos predecibles, pero requieren de muy pocos recursos para su procesamiento y son fáciles de implementar. Un ejemplo clásico de IA basada en funciones hash es el en juego \textit{Space Invaders} \cite{wikipedia_artificial_2025} donde los invasores estaban controlados con funciones hash. 

\subsection{Máquinas de estados finitos}

Las máquinas de estados finitos son un modelo de computación conceptual que describe un sistema que solo puede encontrarse en un estado a la vez, y que puede cambiar a uno de sus otros estados como respuesta a ciertos eventos. Este mecanismo se ha utilizado y sigue utilizándose para un gran número de aplicaciones dentro de los videojuegos. Por ejemplo, en el juego \textit{Pac\-man}, la IA de los fantasmas se rije por el estado en el que se encuentran, como "cazando" o "siendo cazados" \cite{mike_game_2016}. Otro uso está en la gestión de las animaciones de los personajes, donde cada estado corresponde a una posición o movimiento específico de la malla, que al mezclarse mediante interpolaciones generan animaciones. Al igual que las funciones hash, requieren de pocos recursos para su procesamiento, pero su naturaleza determinista y la posibilidad de entrar en bucles sin salida si no están bien diseñadas, limitan su uso en los escenarios más exigentes.

\subsection{Árboles de comportamiento}

Como evolución a las máquinas de estados finitos, surgieron los árboles de comportamiento, permitiendo un mayor grado de complejidad y flexibilidad en la toma de decisiones de los personajes. Los árboles de comportamiento son grafos dirijidos acíclicos, con nodos hoja que representan acciones y varios tipos de control de flujo que determinan el orden de ejecución de las hojas \cite{epic_games_behavior_2025}. Una de sus características más destacadas es su modularidad, lo que permite reutilizar y combinar comportamientos fácilmente. Un fantastico ejemplo de su correcto uso se expuso en la charla de la ``Game Developers Conference'' (GDC) de 2005, sobre la IA de \textit{Halo 2}, donde se explica que su capacidad de reutilización de árboles y el ser más sencillos de entender para los diseñadores del juego, hizo que los programadores optaran por su implementación en lugar de las máquinas de estados finitos \cite{isla_managing_2005}.


\subsection{Planificación de acciones orientada a objetivos}

En contraste con los árboles de comportamiento, que generan una jerarquía de caminos, la planificación de acciones orientada a objetivos (GOAP, por sus siglas en inglés) es más parecida a la planificación de rutas, donde se busca el camino más óptimo para alcanzar un objetivo en función del estado del mundo. Cada uno de los pasos que llevan al objetivo tiene una serie de precondiciones que deben cumplirse, a los cuales se les puede asignar un coste. Es tal su similitud con la planificación de rutas, que este tipo de IA suelen utilizar internamente algoritmos de búsqueda como A*\footnote{A* también se utiliza en la búsqueda de caminos tradicional dentro de los videojuegos, es decir, para saber la ruta que debe tomar un objeto por el mapa para llegar a un destino, y no solo en GOAP.} para encontrar el camino más óptimo. Una de las mejores implementaciones de este tipo de IA se encuentra en el título de 2005 \textit{F.E.A.R.}, donde los enemigos y aliados utilizan GOAP para ``planificar'' su comportamiento una manera mucho más proactiva que sus predecesores \cite{jeff_gdc_2006}.


\subsection{Aprendizaje por refuerzo y redes neuronales}

El aprendizaje por refuerzo (RL, por sus siglas en inglés) es un subcampo del aprendizaje automático que trata de mejorar el comportamiento de un agente a través del feedback que recibe de su entorno. Utilizando un sistema de recompensas y penalizaciones, el agente aprende a tomar decisiones que maximicen su recompensa total a lo largo del tiempo. Aunque el RL empezó con técnicas tabulares, los algoritmos recientes incorporan redes neuronales en diferentes partes del proceso \cite{ghasemi_comprehensive_2025}. Quizás el ejemplo que más se suele ver en la literatura de videojuegos es el uso de RL para controlar vehículos autónomos. Una implementación de esta tecnología en un videojuego comercial está en el título \textit{Star Wars: Outlaws}, donde las motos controladas por IA utilizan un sistema de RL para moverse por el mapa o perseguir al jugador \cite{gaudreau_game_2025}. Sus desarrolladores argumentan que les permitió tener un sistema de control que se adapta a las diferentes rutas del juego sin importar los cambios en el mapa que haya entre versiones.

Pero el uso de redes neuronales no se limita a la toma de decisiones. Su capacidad para realizar inferencias rápidas a partir de datos complejos las hace ideales para otras tareas, como la deformación de mallas. El objetivo en este caso es evitar el fenómeno conocido como el ``valle inquietante'' (\textit{Uncanny Valley}), que se produce cuando una réplica de un ser humano es muy realista, pero no perfecta, generando una sensación de extrañeza en el espectador. Un personaje que sonríe, por ejemplo, pero cuyos músculos faciales alrededor de los ojos permanecen inmóviles, es un caso clásico de este efecto. Para intentar superar este problema, Epic Games optó por crear un sistema para Unreal Engine que utiliza modelos de redes neuronales que simulan el comportamiento de la musculatura y los tejidos blandos bajo la piel al extenderse y contraerse \cite{epic_games_ml_2025}. Un ejemplo de uso de esta reciente tecnología se puede encontrar en la demo técnica que CD Projekt Red mostró en el Unreal Fest de 2025 \cite{cd_projekt_red_witcher_2025}. En ella, se mostraba como se modelaba internamente la musculatura de un caballo para que sus animaciones de movimiento fueran mucho más realistas.

% Me parece que ya es suficiente con lo que hay. Quito esta sección y la siguiente.

% \subsection{Unión de varias técnicas}
% The Last of Us: Parte II --> muchos sistemas entrelazados
% Alien: Isolation --> Árbol de comportamiento + Director de IA

% \subsection{Modelos de lenguaje grande}
% No es una tecnología lo suficientemente madura como para uso actual en videojuegos comerciales.

\section{Metaheurísticas en la investigación científica sobre videojuegos} \label{sec:estado_arte}



% \subsection{Generación de contenido procedural}
% Procedural Generation of Quests for Games Using Genetic Algorithms and Automated Planning --> https://www.icad.puc-rio.br/~logtell/papers/Edirlei_SBGames_2019.pdf
% Search-Based Procedural Content Generation: A Taxonomy and Survey --> https://ieeexplore.ieee.org/document/5756645
% Evaluating Alternative Metaheuristic Algorithms for Procedural Content Generation in Game Design --> https://www.researchgate.net/publication/389270827_Evaluating_Alternative_Metaheuristic_Algorithms_for_Procedural_Content_Generation_in_Game_Design
% Systematic Literature Review of Meta-heuristic Algorithms and their Application in Procedural Content Generation (PCG) in the Context of Computer Games --> https://labs.sciety.org/articles/by?article_doi=10.21203/rs.3.rs-4883187/v1
% Procedural Video Game Scene Generation by Genetic and Neutrosophic WASPAS Algorithms: https://www.mdpi.com/2076-3417/12/2/772
% Flexible Muscle-Based Locomotion for Bipedal Creatures --> https://www.goatstream.com/research/papers/SA2013/
% Comparative Analysis of Metaheuristic Algorithms for Procedural Race Track Generation in Games --> https://www.igi-global.com/article/comparative-analysis-of-metaheuristic-algorithms-for-procedural-race-track-generation-in-games/350330

% Continuous Spatial Public Goods Game Based on Particle Swarm Optimization with Memory Stability --> https://www.mdpi.com/2227-7390/10/23/4572

% AntBot: ant colonies for video games --> https://www.computer.org/csdl/journal/ci/2012/04/06262464/13rRUzpzeDF

% Evolutionary Artificial Intelligence for MOBA / Action-RTS Games using Genetic Algorithms --> https://www.ijcaonline.org/proceedings/icrtitcs2012/number10/10320-1465/#:~:text=The%20Genetic%20Algorithms%20hence%20defines,reactive%20response%20of%20the%20enemy.
% An evolutionary approach to balancing and disrupting real-time strategy games --> https://mssanz.org.au/modsim2021/papers/M8/snell.pdf


\section{Algoritmos evolutivos en juegos de cartas} \label{sec:trabajos_relacionados}

% Hablar sobre el software para la simulación de juegos de cartas: Magic Workstation, Sabberstone, etc.

% No solamente hay simuladores para crear mazos, sino también para crear jugadores virtuales (bots)

% Paper de Pablo y TFG de los ganadores de la competición de IA de Tales of Tribute. Comparar sus enfoques al mio

\subsection{Optimización de agentes de Hearthstone mediante un algoritmo evolutivo}
% Pablo --> https://www.sciencedirect.com/science/article/pii/S0950705119304356

\subsection{Desarrollo de un agente para el concurso de IA Tales of Tribute}
% MCTS es el state of the art en juegos de cartas (buscar como le va al aprendizaje por refuerzo actualmente)

% Meter foto de ganadores del año pasado
\todo{Analizar la estrategia de MCTS de los ganadores de la competición de IA de Tales of Tribute OG7.}
\chapter{Tales of Tribute y Scripts of Tribute}


\section{El videojuego Tales of Tribute}


\section{El motor Scripts of Tribute}


\subsection{La competición ``IEEE Conference on Games''}


\subsection{Partes del motor}

\chapter{Planificación del proyecto} \label{chap:planificacion}

En este capítulo se detallan los aspectos relacionados con los recursos que se han utilizado para el desarrollo del proyecto, así como la metodología que se ha seguido y la distribución temporal utilizada en cada tarea.

\section{Recursos utilizados} \label{sec:recursos_utilizados}

\subsection{Recursos software} \label{sec:recursos_software}

El proyecto se puede dividir en tres partes: el desarrollo del bot en C\#, el ajuste de los pesos en Python y la redacción del informe en LaTeX. Afortunadamente, existe una herramienta de código abierto y en constante mejora en la que se pueden realizar todas estas tareas: Visual Studio Code (VSCode). Este entorno de desarrollo es la navaja suiza que ha estado presente en todo el proceso, permitiendo la edición de código, la gestión de versiones, la redacción del informe en incluso la conexión con interfáz gráfica entre máquinas por SSH. Todo ello gracias a su robusto sistema de extensiones (y la gran comunidad que lo apoya), que permite añadir prácticamente cualquier funcionalidad que se necesite. En el caso de la gestión de versiones, utilizando la integración con Git de VSCode, en todo momento se ha mantenido un repositorio actualizado al día en GitHub, lo que ha permitido la transferencia del progreso entre las dos máquinas utilizadas durante el desarrollo.

Por último, destacar el uso de Zotero como software de gestión de referencias bibliográficas, así como del propio Scripts of Tribute, la base sobre la que se ha erguido el resto del código.

\subsection{Recursos hardware} \label{sec:recursos_hardware}
% Especificaciones de mi PC y el de Pablo

Se han contado con dos ordenadores principales, el primero es el ordenador personal del autor, mientras que el segundo fue proporcionado por el director del proyecto para realizar los experimentos. Estas son sus características:

\begin{itemize}
	\item \textbf{Ordenador personal (Windows 11):}
	      \begin{itemize}
		      \item Procesador: AMD Ryzen 7 5800X
		      \item Memoria RAM: 32 GB DDR5
		      \item Tarjeta gráfica: NVIDIA GeForce RTX 3070
	      \end{itemize}
	\item \textbf{Ordenador de experimentación (Ubuntu 22.04):}
	      \begin{itemize}
		      \item Procesador: Intel i9-12900KF
		      \item Memoria RAM: 128 GB DDR4
		      \item Tarjeta gráfica: NVIDIA GeForce RTX 4060
	      \end{itemize}
\end{itemize}

\section{Metodología de desarrollo} \label{sec:metodologia_desarrollo}
% Metodología ágil mediante sprints (reuniones con Pablo desde el 12 de noviembre de 2024)

\todo{No tiene mucho sentido que esta sección de metodología este fuera de la parte 2, que literalmente se llama Metodología. Revisarlo.}

Para el desarrollo de este Trabajo de Fin de Máster se ha seguido una metodología ágil a base de ``sprints'' regulares. Desde el 12 de noviembre de 2024, se mantuvieron reuniones de forma constante con el director del proyecto, Pablo García Sánchez, aproximadamente cada 3 semanas. En cada una, primero se procedía a repasar el progreso desde la última reunión, a menudo mostrando en tiempo real las nuevas características o los resultados obtenidos. Sobre esta base, el director proponía mejoras o correcciones a realizar para la próxima ocasión. Luego, entre ambos se decidían los siguientes pasos a seguir, mediante una comunicación abierta, lo que permitía ajustar el rumbo del proyecto según las necesidades y los resultados obtenidos.

Aunque las primeras fases del proyecto se realizaron enteramente en el PC personal del autor, en cuanto se necesitó ejecutar el script de ajuste de pesos (o ``entrenador''), se comenzó a utilizar el ordenador proporcionado por el director para ese propósito. Esto permitió realizar experimentos de forma continuada y sencilla, mientras se seguía trabajando en el código del bot en el PC personal. Así se pudo mantener un flujo de trabajo fluido y eficiente, generando código funcional de forma regular y obteniendo resultados cada vez más desarrollados constantemente.


\section{Distribución temporal y cronograma} \label{sec:cronograma}

Para visualizar la distribución temporal de las diferentes tareas realizadas durante el desarrollo de este TFM, se ha elaborado el siguiente diagrama de Gantt (Figura \ref{fig:gantt}).

\begin{figure}[H]
	\centering
	\begin{tikzpicture}[x=1.2cm, y=0.7cm]
		\begin{ganttchart}[
				hgrid,
				vgrid,
				bar/.style={fill=blue!40},
				bar height=0.6,
				bar label font=\footnotesize\color{black!80},
				group right shift=0,
				group top shift=0.6,
				group height=0.3,
				group/.style={fill=black!30},
				title label font=\bfseries\footnotesize,
				title left shift=0,
				title right shift=0,
				title height=1,
				milestone/.style={fill=red, rounded corners=0pt},
				y unit chart=0.6cm,
				x unit=1.2cm
			]{1}{8}
			\gantttitle{2024}{2}
			\gantttitle{2025}{6} \\
			\gantttitle{Nov}{1}
			\gantttitle{Dic}{1}
			\gantttitle{Ene}{1}
			\gantttitle{Feb}{1}
			\gantttitle{Mar}{1}
			\gantttitle{Abr}{1}
			\gantttitle{May}{1}
			\gantttitle{Jun}{1} \\
			\ganttbar{Rev. bibliográfica}{1}{2} \\
			\ganttbar{Estudio de SoT}{2}{4} \\
			\ganttbar{Desarrollo Bot}{2}{5} \\
			\ganttbar{Pruebas Bot}{5}{5} \\
			\ganttbar{Entrenador básico}{4}{5} \\
			\ganttbar{Script gráficas}{5}{6} \\
			\ganttbar{Exp. básicos}{6}{6} \\
			\ganttbar{Entrenador final}{6}{7} \\
			\ganttbar{Exp. finales}{7}{8} \\
			\ganttbar[bar/.style={fill=green!40}]{Escritura TFM}{7}{8} \\
		\end{ganttchart}
	\end{tikzpicture}
	\caption{Distribución temporal de las tareas del proyecto}
	\label{fig:gantt}
\end{figure}

Como se puede observar en el diagrama, la mayoría de tareas se han desarrollado de forma secuencial, aunque algunas se han solapado parcialmente cuando ha sido posible trabajar en paralelo. La escritura del TFM se concentró en el último mes y medio del proyecto, una vez que se contaba con resultados experimentales suficientes para documentar adecuadamente el trabajo realizado. Este es el contenido de cada una de las tareas:
\begin{itemize}
	\item \textbf{Revisión bibliográfica:} se empezó revisando los trabajos ya descritos en la sección \ref{sec:trabajos_relacionados}, además de aprender a jugar a Tales of Tribute y definir el alcance del proyecto.
	\item \textbf{Estudio de Scripts of Tribute:} en esta etapa se estudió el código fuente de Scripts of Tribute, y se empezó a trabajar con C\# como lenguaje de programación, ya que este nunca había sido utilizado por el autor antes.
	\item \textbf{Desarrollo del bot:} se diseñó y desarrolló el bot a base de iteraciones de otros bots existentes, así como del trabajo de Pablo et al. \todo{cita aquí} y otras referencias.
	\item \textbf{Pruebas del bot:} primeras pruebas para comprobar que el funcionamiento del bot era correcto.
	\item \textbf{Entrenador básico:} en un principio, el entrenador solo tenía un modo de funcionamiento básico y estaba peor organizado, pero el ágil desarrollo de esta versión permitió alcanzar resultados rápidamente a modo de prototipo.
	\item \textbf{Script de gráficas:} se creó un script para generar gráficas que visualizan los resultados de los experimentos realizados, así como la función del entrenador que crea los datasets necesarios para su funcionamiento.
	\item \textbf{Experimentos básicos:} experimentos iniciales para evaluar si el rendimiento del bot mejoraba con el ajuste de pesos.
	\item \textbf{Entrenador final:} desarrollo de la versión final del entrenador, incluyendo los diferentes modos de entrenamiento que se describirán más adelante en la sección \ref{sec:mecanismos_evaluacion_fitness}. 
	\item \textbf{Experimentos finales:} experimentos pensados para la evaluación del rendimiento de los distintos modos de entrenamiento del entrenador final.
	\item \textbf{Escritura del TFM:} finalmente se redactó este informe, documentando todo el proceso realizado durante el desarrollo del proyecto.
\end{itemize}

\part{Diseño e implementación: El software desarrollado} \label{part:preparacion}
\chapter{Diseño del agente autónomo} \label{chap:software_bot}

En este capítulo se desglosa el diseño y la implementación del agente autónomo. El enfoque principal se centra en un sistema de toma de decisiones que evalúa cada acción posible mediante una función de puntuación ponderada, cuyos pesos son optimizados por el algoritmo evolutivo descrito en el siguiente capítulo. A continuación, se detallará la arquitectura de este sistema, la definición de los pesos que conforman su ``genoma'' y las lógicas específicas y heurísticas que complementan su comportamiento.

\section{Arquitectura de la toma de decisiones} \label{sec:arquitectura_toma_decisiones}

El comportamiento del agente se basa en una estrategia de decisión voraz con un horizonte de planificación de un solo paso. Cuando el bot debe realizar una acción, el motor del juego le proporciona una lista de todos los movimientos legalmente posibles. Para cada uno de estos movimientos, el bot realiza una simulación interna para averiguar como cambiaría el estado de la partida al aplicarlo. A este nuevo estado se le asigna una puntuación numérica a través de una función de evaluación heurística. Finalmente, el bot selecciona y ejecuta de forma determinista el movimiento que conduce al estado con la puntuación más alta, repitiendo este ciclo hasta que la acción elegida sea la de finalizar el turno. Esto contrasta con el enfoque de MCTS, pues en dicho algoritmo no se simula solo un movimiento, sino la partida entera en múltiples posibles ramificaciones. Una de las ventajas del enfoque voraz es su rapidez, en ningún momento se llegaría ni al límite de tiempo de 10 segundos por turno, ni al de memoria de 256 MB, lo que permite al bot tomar decisiones incluso en situaciones de alta presión temporal. Además, aunque el enfoque MCTS podría ofrecer una mejor exploración de las posibles jugadas, también es cierto que cuanto más se aleja del estado actual de la partida, más incertidumbre hay sobre el valor de las jugadas debido a la alta estocasticidad del juego y menos valor tendría la evaluación de un estado futuro lejano.

El punto más importante este sistema es la función de evaluación, ``ScoreMove'', que calcula la idoneidad de un estado del juego. Esta función es la que ejecuta la evaluación lineal ponderada, donde se calculan las diferencias entre el estado del juego antes y después de la jugada simulada a través de un conjunto de características. Cada una de estas características se multiplica por un peso específico, y la suma (o resta, en caso de gestionar datos relativos al contrincante) de todos estos componentes da como resultado la puntuación final de la acción.

La fórmula general que representa este cálculo se puede expresar según la ecuación \ref{eq:puntuacion_movimiento}:

\begin{equation}
\text{Puntuación}(m) = \sum_{i=1}^{n} w_i \cdot \Delta f_i(E_{t+1}, E_t) \label{eq:puntuacion_movimiento}
\end{equation}

Donde $\text{Puntuación}(m)$ es el valor asignado al movimiento $m$, $n$ es el número total de características evaluadas, $w_i$ es el peso asociado a la característica $i$, y $\Delta f_i$ es la función que calcula la diferencia de valor de dicha característica entre el estado resultante $E_{t+1}$ y el estado original $E_t$. Para garantizar que la comparación entre los diferentes movimientos sea justa y reproducible, todas las simulaciones internas se realizan utilizando una semilla pseudoaleatoria constante, lo que elimina la estocasticidad de ciertos eventos del juego, como el robo de cartas de la taberna. En el siguiente pseudocódigo se muestra el bucle de evaluación que sigue el bot para cada uno de los posibles movimientos:

\begin{algorithm}[H]
    \caption{Selección de la Mejor Acción en un Turno}
    \label{alg:get_best_move}
    \begin{algorithmic}[1]
        \State \Comment{Lógica interna del agente para elegir una acción en el estado de juego actual $S$.}
        \State \textbf{Entradas:} Estado del juego $S$, Lista de acciones posibles $A = \{a_1, ..., a_n\}$, Pesos $\vec{w}$.
        \State
        \State $\text{mejorAcción} \leftarrow \text{AcciónAleatoria}(A)$
        \State $\text{mejorPuntuación} \leftarrow -\infty$
        \State
        \State \textbf{para cada} acción $a_i \in A$ \textbf{hacer}
        \State \quad \textbf{si} $a_i$ es ``Pasar Turno'' \textbf{entonces}
        \State \qquad $\text{puntuación} \leftarrow -1$
        \State \quad \textbf{sino}
        \State \qquad $S_{t+1} \leftarrow \text{Simular}(S, a_i)$
        \State \qquad $\text{puntuación} \leftarrow \text{Evaluar}(S_{t+1}, S, \vec{w})$ \Comment{Usa la Ecuación \ref{eq:puntuacion_movimiento}.}
        \State \quad \textbf{fin si}
        \State
        \State \quad \textbf{si} $\text{puntuación} > \text{mejorPuntuación}$ \textbf{entonces}
        \State \qquad $\text{mejorPuntuación} \leftarrow \text{puntuación}$
        \State \qquad $\text{mejorAcción} \leftarrow a_i$
        \State \quad \textbf{fin si}
        \State \textbf{fin para}
        \State
        \State \textbf{devolver} $\text{mejorAcción}$
    \end{algorithmic}
\end{algorithm}

\section{Definición de los pesos} \label{sec:definicion_pesos}

El comportamiento estratégico del bot está enteramente definido por un conjunto de 20 pesos numéricos. Este vector de valores representa el ``genoma'' de un individuo dentro del algoritmo evolutivo. Cada peso está asociado a una característica específica del juego, permitiendo al proceso de optimización ajustar con precisión la forma en que el bot valora cada aspecto de la partida. La Tabla \ref{tab:pesos_ebw} detalla cada uno de estos pesos y su función dentro de la evaluación.

La transferencia de estos pesos desde el entrenador al bot es un componente clave del sistema. Para evitar la sobrecarga de reescribir ficheros constantemente, el mecanismo principal se basa en el uso de variables de entorno. El script de Python convierte el genoma de un individuo (un array de 20 números) en una cadena de texto separada por comas. Esta cadena se asigna a una variable de entorno antes de lanzar el proceso del ``GameRunner''. Al iniciarse, el bot detecta estas variables, las lee, y utiliza sus valores para configurar la función de evaluación para esa partida en concreto.

	{\scriptsize
		\setlength{\tabcolsep}{4pt}
		\begin{longtable}{@{}l>{\tiny\raggedright\arraybackslash}p{2.2cm}p{0.65\textwidth}@{}}
			\caption{Descripción de los pesos del genoma del agente evolutivo.} \label{tab:pesos_ebw}                                                                                       \\

			\toprule
			\textbf{Categoría} & \textbf{Peso (Enum)} & \textbf{Descripción}                                                                                                                \\
			\midrule
			\endfirsthead

			\multicolumn{3}{c}{{\tablename\ \thetable{} -- Continuación}}                                                                                                                   \\
			\toprule
			\textbf{Categoría} & \textbf{Peso (Enum)} & \textbf{Descripción}                                                                                                                \\
			\midrule
			\endhead

			\midrule
			\multicolumn{3}{r}{{Continúa en la siguiente página}}                                                                                                                           \\
			\endfoot

			\bottomrule
			\endlastfoot

			\multirow{4}{*}{\begin{tabular}[c]{@{}l@{}}Agentes y\\ Combate\end{tabular}}
			                   & A\_HEALTH\_ REDUCED  & Penalización por el daño recibido por un agente propio, o bonificación por el daño infligido a un agente enemigo.                   \\
			                   & A\_KILLED            & Penalización si un agente propio es destruido, o bonificación equivalente si se destruye un agente enemigo.                         \\
			                   & A\_OWN\_AMOUNT       & Pondera la importancia de tener un mayor número de agentes en el tablero.                                                           \\
			                   & A\_ENEMY\_ AMOUNT    & Pondera la importancia de que el enemigo tenga un mayor número de agentes en el tablero.                                            \\
			\midrule
			\multirow{8}{*}{\begin{tabular}[c]{@{}l@{}}Recursos y\\ Cartas\end{tabular}}
			                   & CURSE\_REMOVED       & Bonificación por eliminar una carta de maldición del propio mazo.                                                                   \\
			                   & C\_TIER\_POOL        & Valor asignado a las cartas del mazo según una clasificación de tier predefinida.                                                   \\
			                   & C\_TIER\_TAVERN      & Penalización por dejar cartas de tier alto disponibles para el oponente en la taberna.                                              \\
			                   & C\_GOLD\_COST        & Valor asignado a las cartas en función de su coste en oro.                                                                          \\
			                   & C\_OWN\_COMBO        & Bonificación por tener varias cartas del mismo patrón en el mazo, aumentando el potencial de combo.                                 \\
			                   & C\_ENEMY\_COMBO      & Pondera el potencial de combo del mazo del oponente.                                                                                \\
			                   & COIN\_AMOUNT         & Pondera la diferencia de monedas generadas en un turno.                                                                             \\
			                   & POWER\_AMOUNT        & Pondera la diferencia de poder generado en un turno.                                                                                \\
			                   & PRESTIGE\_AMOUNT     & Pondera la diferencia de prestigio ganado o perdido.                                                                                \\
			                   & H\_DRAFT             & Peso que se utiliza como multiplicador defensivo para valorar la acción de quitarle una carta útil al oponente (``hate drafting''). \\
			\midrule
			\multirow{5}{*}{\begin{tabular}[c]{@{}l@{}}Cartas\\ Específicas\end{tabular}}
			                   & T\_TITHE             & Valor asignado a la carta ``Tithe'', que permite una activación de patrón adicional.                                                \\
			                   & T\_BLACK\_ SACRAMENT & Valor asignado a ``Black Sacrament'', que permite destruir un agente enemigo.                                                       \\
			                   & T\_AMBUSH            & Valor asignado a ``Ambush'', que permite destruir hasta dos agentes enemigos.                                                       \\
			                   & T\_BLACKMAIL         & Valor asignado a ``Blackmail'', que proporciona poder o prestigio.                                                                  \\
			                   & T\_IMPRISONMENT      & Valor asignado a ``Imprisonment'', similar a ``Blackmail'' pero con mayor potencial.                                                \\
			\midrule
			\multirow{1}{*}{Patrones}
			                   & P\_AMOUNT            & Valor asignado a cualquier acción que cambie el favor de un patrón.                                                                 \\
		\end{longtable}
	}

\section{Lógicas específicas y heurísticas} \label{sec:logicas_especificas}

Para aquellos casos especiales en los que el bot puede no ser capaz de discernir la importancia de una jugada, como es el caso de la aparición en la taberna de ciertas cartas importantes, se implementaron una serie de lógicas específicas y heurísticas codificadas manualmente pero con su propio peso asignado.

El ejemplo más representativo se encuentra en la función ``CalculateTavernScore''. Esta función no evalúa el valor de adquirir una carta de la taberna para uno mismo, sino que calcula el valor estratégico de eliminarla de la taberna para que el oponente no pueda adquirirla. Esta técnica, popularmente denominada como ``hate drafting'', se implementó tras ver a algunos jugadores explicar sus movimientos mientras jugaban. En los vídeos se observaba que a veces preferían quitarle una carta al oponente aunque no les fuera útil a ellos mismos, pues que el rival la utilizase sería peor que los recursos gastados en comprar la carta. Este cálculo es puramente heurístico y depende del estado actual de la partida. Por ejemplo, para la carta ``Tithe'', que permite una activación de patrón adicional, el bot la valora más positivamente de forma defensiva (es decir, quitarla para que el oponente no la compre) si el oponente está cerca de alcanzar la victoria por prestigio, ya que un turno con dos acciones de patrón podría ser decisivo para él. De forma similar, se evalúan otras cartas clave como ``Ambush'', cuyo valor aumenta exponencialmente cuantos más agentes tenga el oponente en el tablero.

Además de esta lógica centrada en la taberna, existen otras heurísticas integradas en la evaluación. En ``CalculateScoreAgents'', por ejemplo, se diferencia entre simplemente dañar a un agente y destruirlo por completo, asignando una penalización mucho mayor a esto último. Asimismo, en la función ``ComputeCardPoolValue'' se aplica una lógica especial para las cartas de ``maldición'': tenerlas en el propio mazo es negativo, pero conseguir que el oponente adquiera una es positivo. Esta combinación de una función de evaluación general y optimizable con heurísticas específicas para casos concretos intenta unificar el conocimiento del jugador con la potencia bruta de la optimización evolutiva. El sistema evolutivo se encarga de ajustar el comportamiento general, mientras que el conocimiento codificado en las heurísticas garantiza que el bot actúe de forma decisiva en momentos clave de la partida.

El desarrollo del bot marca el cumplimiento del \textbf{OG2}, que establece el diseño e implementación un agente inteligente en C\# para ``Scripts of Tribute'', cuya toma de decisiones se base en una evaluación heurística del estado del juego mediante una función de fitness ponderada, incorporando conocimiento específico del dominio.
\chapter{Sistema de entrenamiento evolutivo}
% La parte en Python

\section{Arquitectura general del entrenador}
% Visión general del sistema. Script en Python que usa Inspyred para orquestar un algoritmo evolutivo
% Encontrar el conjunto de pesos óptimo para el bot
% Diagrama de flujo

\section{Comunicación entrenador-imulador}
% Los subprocesos y las variables de entorno
% os.environ para pasar los pesos de un individuo al bot en C#

\section{El algoritmo evolutivo con Inspyred}
% Justificar la elección de Inspyred
% Describir los componentes de un EA: individuos, población, generación inicial, operadores de variación, evaluación del fitness, selección, etc
% Explicar los valores por defecto de Inspyred

\section{Paralelización del algoritmo}
% Uso de multiprocessing para el diccionario y concurrent.futures para el evaluador

\section{Mecanismos de evaluación del fitness}

\subsection{Modo fijo: evaluación contra oponentes estáticos}
% El individuo juega contra un conjunto predefinido de bots. El fitness es el número de victorias
% inspyred.ec.evaluators.parallel_evaluation_mp
\subsection{Modo coevolución: competición interna}
% Los individuos (padres e hijos) compiten entre sí en un torneo round robin. El fitness es el número de victorias
% evaluate_coevolution_orchestrator
\subsection{Modo híbrido: combinando estrategias}
% Sistema de entrenamiento por segmentos hybrid_schedule_str
% Gestión de cambios de modo
\section{El salón de la fama}
% Mantener un diccionario de oponentes de élite (catastrophic forgetting)

\part{Experimentación y conclusiones} \label{part:resultados}
\chapter{Diseño experimental} \label{chap:experimentacion}

Esta tercera parte del TFM, aborda el diseño experimental, los resultados y las conclusiones del proyecto. Una vez descritos los componentes individuales, es decir, el agente autónomo y el sistema de entrenamiento evolutivo, este capítulo se centra en el diseño de los experimentos llevados a cabo para evaluar y comparar las distintas estrategias de optimización. El objetivo de la fase experimental no es solo encontrar el mejor conjunto de pesos posible, sino también entender cómo las diferentes configuraciones del algoritmo evolutivo afectan al proceso de aprendizaje y a la naturaleza de las soluciones encontradas. Para ello, se han diseñado una serie de experimentos que permiten analizar el impacto de factores como el modo de evaluación (fijo, coevolutivo o híbrido) y el uso de un salón de la fama.

\section{Configuración de los experimentos} \label{sec:configuracion_experimentos}

El diseño de los experimentos se ha centrado en responder a tres preguntas fundamentales: ¿qué combinación de estrategias en el modo híbrido produce los mejores agentes?, ¿qué modo de entrenamiento genera los mejores bots? y ¿cuál es el impacto del salón de la fama en el rendimiento del entrenamiento? Para abordar estas cuestiones, se crearon scripts de ejecución que automatizan el lanzamiento de varios conjuntos entrenamientos, cada uno con una configuración específica. Cada experimento se conforma de una serie de entrenamientos, los cuales a su vez realizan 5 carreras diferentes cada uno. Los resultados de estas carreras se promedian para obtener una medida más estadísticamente robusta del rendimiento de cada entrenamiento.


\subsection{Análisis de configuraciones híbridas} \label{sec:experimentos_hibridos}

Para explorar en profundidad el modo híbrido, se diseñó un experimento dedicado exclusivamente a probar una veintena de configuraciones distintas de este modo. Sus configuraciones varían tanto en el número de segmentos como en la proporción de evaluaciones dedicadas a cada uno de ellos, con el objetivo de encontrar un balance óptimo entre la explotación de conocimiento (evaluación contra oponentes fijos) y la exploración de nuevas estrategias (competición interna). Dado que en este experimento se busca comparar únicamente el efecto de la variación en el parámetro que controla el orden y proporción de los segmentos que conforman el modo híbrido, se mantuvieron constantes los demás parámetros del entrenador. Esta decisión permite aislar las diferencias de rendimiento únicamente al impacto de dicho parámetro, pero también tiene un efecto colateral: el tiempo de entrenamiento de cada configuración es diferente. Dado que el modo híbrido cuenta con diferentes proporciones de evaluaciones fijas y coevolutivas, el tiempo de entrenamiento puede variar significativamente entre cada configuración. Sin embargo, este aspecto se puede considerar como algo positivo, ya que también permite observar el rendimiento de cada configuración en función del tiempo de entrenamiento, pudiendo encontrar configuraciones que no solo sean más efectivas, sino que también sean más eficientes en términos temporales.

La Tabla \ref{tab:hybrid_schedules} detalla cada una de las 10 configuraciones que se evaluaron. Para cada una, se ha calculado el número total de segmentos, el porcentaje de evaluaciones dedicadas a cada modo (fijo para la explotación y coevolución para la exploración), y una breve descripción de la estrategia que representa. Los parámetros comunes del experimento fueron: 5 carreras de entrenamiento por experimento, una población de 10 individuos, 500 evaluaciones máximas, 200 partidas por evaluación, 8 hilos de paralelismo, y la desactivación tanto del salón de la fama global como del interno para aislar el efecto de la planificación híbrida. Como se han asignado 500 evaluaciones a cada entrenamiento, y cada carrera consta de 10 rondas, esto implica que cada entrenamiento se ejecuta durante un total de 50 generaciones.

\small
\begin{longtable}{@{}l p{7cm} c c c@{}}

	\caption{Configuraciones de entrenamiento en modo híbrido evaluadas.}
	\label{tab:hybrid_schedules}                                                                                                 \\
	\toprule
	\textbf{ID} & \textbf{Planificación (\texttt{hybrid\_schedule\_str})} & \textbf{Seg.} & \textbf{\% Fijo} & \textbf{\% Coevo} \\
	\midrule
	\endfirsthead

	\multicolumn{5}{c}%
	{{\bfseries -- continuación de la Tabla \thetable --}}                                                                       \\
	\toprule
	\textbf{ID} & \textbf{Planificación (\texttt{hybrid\_schedule\_str})} & \textbf{Seg.} & \textbf{\% Fijo} & \textbf{\% Coevo} \\
	\midrule
	\endhead

	\bottomrule
	\endfoot

	H-1         & \texttt{fixed:0.4,coevolution:0.3,fixed:0.3}            & 3             & 70\%             & 30\%              \\
	H-2         & \texttt{fixed:0.7,coevolution:0.3}                      & 2             & 70\%             & 30\%              \\
	H-3         & \texttt{fixed:0.4,coevolution:0.2,fixed:0.4}            & 3             & 80\%             & 20\%              \\
	H-4         & \texttt{fixed:0.5,coevolution:0.5}                      & 2             & 50\%             & 50\%              \\
	H-5         & \texttt{coevolution:0.6,fixed:0.4}                      & 2             & 40\%             & 60\%              \\
	H-6         & \texttt{coevolution:0.4,fixed:0.2, coevolution:0.4}     & 3             & 20\%             & 80\%              \\
	H-7         & \texttt{fixed:0.8,coevolution:0.2}                      & 2             & 80\%             & 20\%              \\
	H-8         & \texttt{coevolution:0.8,fixed:0.2}                      & 2             & 20\%             & 80\%              \\
	H-9         & \texttt{fixed:0.3,coevolution:0.4,fixed:0.3}            & 3             & 60\%             & 40\%              \\
	H-10        & \texttt{coevolution:0.3,fixed:0.4, coevolution:0.3}     & 3             & 40\%             & 60\%              \\
\end{longtable}

\subsection{Análisis de los modos de entrenamiento y el salón de la fama} \label{sec:experimentos_modos_entrenamiento}

Para hallar el método que genera los bots con mayor porcentaje de victoria se creó un experimento que ejecuta los tres modos de entrenamiento principales (fijo, coevolutivo e híbrido) dos veces. La primera vez con el salón de la fama activado (con un tamaño de 3 individuos), y la segunda vez completamente desactivado. Por lo tanto, este segundo experimento trata de aportar información en dos frentes, tanto a nivel de modo de entrenamiento como a nivel de salón de la fama. El problema del tiempo de entrenamiento persiste en este experimento, pero acentuado aun más, pues los parámetros entre cada configuración pueden variar significativamente, no solo en el número de evaluaciones, sino también en el número de miembros de la población. Dado que este experimento no trata de comparar un único parámetro, sino configuraciones completamente diferentes, en este caso sí se optó por usar el tiempo como medida de unificación. Así, se ha tratado que las configuraciones se ejecuten durante un tiempo similar, permitiendo una comparación más justa entre las distintas estrategias de entrenamiento.

La Tabla \ref{tab:main_experiments} resume las 6 configuraciones específicas que se ejecutaron. Se puede observar cómo los parámetros de ``Tamaño de Población'' y ``Evaluaciones Máximas'' se ajustaron para cada modo. Para todos los experimentos de esta serie se mantuvieron constantes los siguientes parámetros base: 5 carreras de entrenamiento por experimento, 200 partidas por evaluación y 8 hilos de paralelismo para los procesos de Python.

\todo{Revisar parámetros finales}

\begin{table}[H]
	\centering
	\caption{Configuraciones para la comparativa de modos de entrenamiento y salón de la fama.}
	\label{tab:main_experiments}
	\begin{tabular}{@{}llcccl@{}}
		\toprule
		\textbf{ID} & \textbf{Modo} & \textbf{Población} & \textbf{Eval.} & \textbf{Salón F.} & \textbf{Plan híbrido}      \\ \midrule
		E-1         & Fijo          & 10                 & 500            & 3                 & N/A                        \\
		E-2         & Híbrido       & 10                 & 500            & 3                 & \texttt{f:0.4,c:0.2,f:0.4} \\
		E-3         & Coevolución   & 10                 & 500            & 3                 & N/A                        \\
		\midrule
		E-4         & Fijo          & 10                 & 500            & 0                 & N/A                        \\
		E-5         & Híbrido       & 10                 & 500            & 0                 & \texttt{f:0.4,c:0.2,f:0.4} \\
		E-6         & Coevolución   & 10                 & 500            & 0                 & N/A                        \\ \bottomrule
	\end{tabular}
\end{table}

\section{Métricas de evaluación} \label{sec:metricas_evaluacion}

Para analizar los resultados de estos experimentos de una forma cuantitativa y visual, se desarrolló un script en Python que procesa los ficheros de datos en formato CSV generados por el entrenador. Este script genera un conjunto de gráficas estandarizadas que permiten evaluar y comparar el rendimiento de las diferentes configuraciones a través de varias métricas clave.

Una de las métricas principales es la evolución del \textit{fitness} a lo largo de las generaciones. Se genera una gráfica de líneas que muestra el \textit{fitness} máximo, medio y mínimo, promediado a través de todas las carreras de entrenamiento de un mismo experimento. Esta visualización pretende utilizarse para entender la velocidad de convergencia y si el algoritmo se estanca en óptimos locales.

Otra área de análisis se centra en la convergencia y distribución de los pesos del genoma. Para ello se utilizan dos tipos de gráficas. Primero, un mapa de calor que muestra la evolución del valor promedio de cada uno de los 20 pesos a lo largo de las generaciones. Segundo, un diagrama de barras que presenta el valor medio final de cada peso en la última generación. Estas gráficas permiten analizar si el algoritmo converge hacia soluciones ``generalistas'', donde muchos pesos tienen valores moderados, o hacia soluciones ``especialistas'', donde unos pocos pesos tienden a valores extremos (cercanos a 0 o 1), indicando que el bot ha aprendido a priorizar unas pocas heurísticas muy específicas. Para medir la consistencia y el rendimiento del campeón final, se utilizan dos enfoques. Por un lado, un diagrama de cajas para cada peso muestra la distribución de los valores finales de los campeones de todas las carreras. El script permite generar todas estas gráficas tanto de los campeones de cada generación como de la media de los pesos de los individuos de cada generación, lo que proporciona una visión desde varios ángulos del proceso de entrenamiento.

Finalmente, la métrica de rendimiento definitiva es el benchmark que se ejecuta al concluir cada entrenamiento. En esta fase, los campeones de cada carrera compiten entre sí para determinar al ganador absoluto, cuya tasa de victorias en este torneo final se considera la medida última del éxito de esa configuración de entrenamiento. Además, de entre los campeones de cada entrenamiento (que ya son los campeones de su conjunto de carreras), se realiza un segundo torneo final para determinar el campeón absoluto de ese experimento. Este torneo final se ejecuta contra un conjunto de bots fijos, que actúan como una referencia estable para evaluar el rendimiento del agente entrenado en un contexto parecido al de la competición real CoG.

Este conjunto de scripts para la generación de gráficas y benchmarking de agentes, cuyos resultados se utilizarán en el siguiente capítulo, demuestran el cumplimiento del objetivo específico \textbf{OG4} del proyecto, que se centra en el desarrollo de herramientas para la visualización y análisis de los datos generados durante el entrenamiento.

\subsection{Bots de referencia} \label{sec:bots_referencia}

Durante toda esta memoria se ha hecho alusión a los agentes autónomos que se han usado en el modo fijo y en los torneos finales para evaluar el rendimiento de los agentes evolutivos. Estos bots de referencia vienen incluidos en \textit{Scripts of Tribute-Core} y utilizan diferentes algoritmos y heurísticas para cumplir su función. A día de la redacción de este documento, existen 7 bots de referencia más el bot \textit{Random}, que actúa como un agente completamente aleatorio. Lastimosamente, no todos funcionan sin errores. La razón de estos problemas de ejecución es la misma por la que no se han podido utilizar los bots ganadores de competiciones anteriores durante el entrenamiento evolutivo: el motor de simulación SoT se ha actualizado con nuevas funciones y cambios en la API, lo que ha provocado que algunos bots antiguos no sean compatibles con la versión actual del motor. Se contactó con los creadores de SoT para intentar que solucionaran estos problemas de incompatibilidad y, si bien están dispuestos a arreglar los bots en algún momento, la prioridad para ellos es el motor en sí. Por lo tanto, los agentes que sí se han podido utilizar son:

\begin{itemize}
	\item \textbf{PatronFavorsBot}: bot simple que se centra en ganar el favor de los patrones a toda costa. En cada movimiento comprueba si puede ganar el favor de un patrón y, si no es así, realiza un movimiento aleatorio.
	\item \textbf{MaxAgentsBot}: este también es un bot simple, en este caso se centra en maximizar el número de agentes que tiene en el tablero. En cada movimiento, comprueba si puede comprar o colocar un agente, y si no es así, realiza un movimiento aleatorio. Además, para elegir entre los bots a comprar o colocar utiliza una heurística que prioriza las cartas de una mejor calidad.
	\item \textbf{MaxPrestigeBot}: similar a los bots anteriores, este bot se centra en maximizar su puntuación de prestigio. Para cada acción posible, evalúa cual es la que le otorga más puntos de prestigio y la ejecuta. Si no puede realizar ninguna acción que le otorgue puntos, realiza un movimiento aleatorio.
	\item \textbf{DecisionTreeBot}: este es el único bot que no se basa en una heurística simple, sino que utiliza un árbol de decisión diseñado por un experto. Es decir, es un bot con un esquema de funcionamiento completamente opuesto al \textit{EvolutionaryBot} de este proyecto, ya que está diseñado para tomar decisiones basadas en un conjunto de reglas fijas y no evoluciona ni aprende de la experiencia. En lugar de evaluar y puntuar todos los movimientos posibles de forma unificada, el bot sigue una secuencia predefinida: primero intenta jugar las cartas de su mano, si no es posible, considera comprar cartas de la taberna, después activar un patrón, y así sucesivamente hasta encontrar una acción válida. La complejidad de este agente no reside en una función de evaluación ponderada, sino en las heurísticas específicas que gobiernan cada tipo de decisión. Por ejemplo, para decidir qué carta jugar, no elige simplemente la que da más recursos, sino que utiliza una lógica sofisticada que clasifica las cartas en mano según su potencial de combo. En ocasiones, ``sacrifica'' una carta de bajo potencial para preparar el terreno y maximizar los efectos de combo de las cartas que jugará a continuación. De forma similar, para la compra de cartas, utiliza una función heurística que valora la sinergia (favoreciendo cartas de patrones que ya posee) y la calidad general de la carta. Durante la primera fase del proyecto se estudió este bot con detalle, pues su autor claramente tenía un conocimiento profundo del juego y de las estrategias óptimas (algo que no se puede decir del autor de esta memoria). El resultado de esos estudios fueron las heurísticas específicas que el \textit{EvolutionaryBot} utiliza para evaluar la compra y uso de cartas específicas.
\end{itemize}

En cada evaluación fija, el número máximo de partidas configurado se reparte entre estos cuatro bots de referencia, ocurriendo la mitad para cada uno como P1 y la otra mitad como P2. Si por ejemplo el número fuera 200, en la evaluación del \textit{fitness} de un individuo se jugarían 50 partidas contra cada uno de los bots de referencia, 25 como P1 y 25 como P2.
\chapter{Resultados}

\section{A nivel de población}


\subsection{Evolución del fitness}


\subsection{Pesos medios}


\subsection{Pesos a lo largo de las generaciones}


\subsection{Variación de los pesos finales}


\section{A nivel de líderes}


\subsection{Evolución del fitness}


\subsection{Pesos medios}


\subsection{Pesos a lo largo de las generaciones}


\subsection{Variación de los pesos finales}


\section{Más resultados...}
\chapter{Conclusiones} \label{chap:conclusiones}




\section{Líneas de trabajo futuro} \label{sec:trabajo_futuro}
% Expandir el conjunto de oponentes fijos con los ganadores de la competición CoG.
% Añadir más heurísticas, como quitarse cartas malas usando la tesorería.
% Optimizar la evaluación con el uso de cartas instantáneas.
% Hyper parameter tuning de los parámetros del algoritmo genético de inspyred: Probar operadores de mutación y cruce más sofisticados.
% Mejorar la elección de mazos mediante winrates.
% Reimplementar el agente y entrenador en un único lenguaje de programación.

La unión de heurísticas específicas junto con los pesos generados por el algoritmo evolutivo han permitido crear un bot con rendimiento muy competitivo, al menos contra los agentes autónomos de \textit{Scripts of Tribute}. Sin embargo, aun existen muchas áreas de mejora que pueden ser exploradas en el futuro, por ejemplo, endurecer esos mismos oponentes durante el entrenamiento podría desembocar en pesos que produzcan dinámicas más capaces de adaptarse a diferentes estilos de juego. El camino más claro para obtener mejores oponentes es conseguir que los participantes de las competiciones de otros años funcionen correctamente en la versión actual del motor. Lo cual no sería nada sencillo, pues cada uno de esos bots son más complejos que los que SoT proporciona por defecto, pero la recompensa podría ser un gran salto cualitativo en el rendimiento del agente. Un aspecto secundario positivo de este enfoque es que se podrían revisar las heurísticas que dichos bots utilizan, las cuales podrían reimplementarse en el agente de este proyecto y así incluir más conocimiento experto en el mismo.

Aunque en este trabajo se ha puesto especial atención a qué oponentes se enfrentan al agente en los enfrentamientos, así como el orden y la cantidad de partidas que estos juegan, no se ha explorado en profundidad el impacto de alterar los algoritmos base que conforman la estrategia evolutiva. En este sentido, se podrían a utilizar operadores de mutación y cruce más sofisticados de entre los disponibles en la librería \textit{inspyred}, incluso llegando a implementar un ajuste de hiperparámetros para encontrar los que maximicen la exploración y la explotación en diferentes etapas del entrenamiento de forma análoga al modo híbrido de entrenamiento.

Otra forma de mejorar el rendimiento del agente es sabiendo qué mazos son más efectivos al ser jugados por el agente. Durante todo este proyecto se ha utilizado un selector aleatorio de mazos, lo cual es lo más indicado para mejorar la exploración durante el entrenamiento, pero una vez se tiene el conjunto de pesos final, es el momento de afinar la potencia existente lo máximo posible. Para ello, simplemente se tendrían que analizar un gran número de partidas y ver cuáles son los mazos que más victorias obtienen. Con esta información, se podría implementar un selector más sofisticado que elija el mazo con mayor probabilidad de ganar en cada partida.

Por último, un cambio a un nivel más holístico sería la refactorización del agente y el entrenador en un único lenguaje de programación. Aunque la idea de utilizar variables de entorno y lecturas de la salida estándar es una técnica perfectamente válida para la versión actual del proyecto, es posible que la unificación de ambos componentes bajo Python permita mejorar el rendimiento del entrenamiento. Para ello, se podría utilizar la implementación del método de comunicación ``gRPC'' que se añadió a SoT recientemente, el cual permitiría que toda la lógica del agente y el entrenador se ejecutara bajo un único programa orquestador, dejando al GameRunner únicamente como un componente de entrada y salida de datos mediante microservicios. Este enfoque permitiría que el entrenador y el agente compartieran objetos y memoria directamente, y además abriría la puerta al uso de otras técnicas más avanzadas como el aprendizaje por refuerzo profundo, que requiere de una alta transferencia de datos del entorno entre el agente y el GameRunner.

% Material de referencia
\backmatter

% % Bibliografía
\cleardoublepage
\bibliographystyle{plainurl}
\bibliography{bibliografia}

% \appendix
% \input{apendices/codigo}

\end{document}